\documentclass[11pt]{article}
\usepackage{fontspec}
\setmainfont{Times New Roman}
\setsansfont{Arial}
\setmonofont{Consolas}
\usepackage{amsmath,amssymb,mathtools}
\usepackage{geometry}
\usepackage{microtype}
\geometry{a4paper,margin=2.28cm}
\setlength{\parindent}{1.15em}
\setlength{\parskip}{0pt}
\makeatletter
\renewcommand\footnotesize{\@setfontsize\footnotesize{7.5}{8.5}\selectfont}
\makeatother
\providecommand{\mG}{\mathfrak{G}}
\providecommand{\mA}{\mathfrak{A}}
\providecommand{\mB}{\mathfrak{B}}
\providecommand{\mE}{\mathfrak{E}}
\providecommand{\mH}{\mathfrak{H}}
\providecommand{\mC}{\mathfrak{C}}
\providecommand{\mD}{\mathfrak{D}}
\providecommand{\mR}{\mathfrak{R}}
\providecommand{\mS}{\mathfrak{S}}
\providecommand{\mT}{\mathfrak{T}}
\providecommand{\mo}{\mathfrak{o}}
\providecommand{\iso}{\simeq}
\emergencystretch=120em
\allowdisplaybreaks
\sloppy
\hbadness=10000
\vbadness=10000
\hfuzz=2pt
\vfuzz=10pt
\raggedbottom
\begin{document}
% Унит: Папер34 сецтион04 в001. Дирецт интерславиц транслатион фром аудитед Герман соурце линес 184--253, сегментс P34-С0051--P34-С0065. Соурце сцан пагес 10--12 ретаинед ас витнесс, and Зенодо рецорд 20822156 вас цхецкед ат 2026-06-24Т07:42З.
\setcounter{footnote}{10}

\subsection*{§ 4. Директне продукты и пресекы}

Група $\mG$ называје се директным продуктом двух факторов, $\mG=\mA\times\mB$, ако
\[
\begin{array}{ll}
1.& \mA,\mB\text{ сут нормалне делительи в }\mG,\\
2.& \mA\mB=\mG,\\
3.& \mA\cap\mB=\mE.
\end{array}
\]
\emph{Дефиниција јест очивидно еквивалентна с наступным:} всакы елемент $g$ од $\mG$ једнозначно представльаје се како $g=ab$, кде $a$ јест в $\mA$, $b$ јест в $\mB$, и елементы од $\mA$ комутујут с елементами од $\mB$: $ab=ba$.

Група $\mG$ называје се директным продуктом $n$ факторов, $\mG=\mA_1\times\cdots\times\mA_n$, ако, положивши $\mB_i=\mA_1\cdots\mA_{i-1}\mA_{i+1}\cdots\mA_n$, имајемо, же $\mG=\mA_i\times\mB_i$ јест директны продукт за всако $i$.

\emph{Дефиниција знову јест еквивалентна с наступным:} всакы $g$ од $\mG$ једнозначно представльаје се како
\[
        g=a_1a_2\cdots a_n,\qquad a_i\in\mA_i,
\]
и елементы од $\mA_i$ комутујут с елементами од $\mA_k$.

Наступне факты често се користајут:
\begin{enumerate}
\item Ако $\mA_1\times\cdots\times\mA_n=\mR$ и $\mR\times\mH=\mG$, тогда $\mA_1\times\cdots\times\mA_n\times\mH=\mG$.
\item Ако $\mG=\mA_1\times\cdots\times\mA_n=\mC_1\times\cdots\times\mC_n$ и $\mC_i\subseteq\mA_i$, тогда $\mC_i=\mA_i$. (То следује из представјеньа елементов $\mA_i$ посредством $\mC_j$.)
\item Ако $\mG=\mA\times\mB$ и $\mR\supseteq\mA$, тогда $\mR=\mA\times(\mR\cap\mB)$. (То следује из представјеньа елементов $\mR$ в форме $ab$, кде вторы фактор належи и до $\mR$, и до $\mB$.)
\item Ако $\mG=\mA\times\mB$, тогда $\mA\sim\mG/\mB$ и $\mB\sim\mG/\mA$ (втораја теорема изоморфизма). Из того далье следује:
\item Ако $\mG=\mA\times\mB$, и $\mA$ и $\mB$ имајут композицијне рјады долгости $m$ и $n$, тогда $\mG$ имаје композицијны рјад долгости $m+n$.
\item Ако $\mA\sim\overline{\mA}$ и $\mB\sim\overline{\mB}$, тогда $\mA\times\mB\sim\overline{\mA}\times\overline{\mB}$.
\end{enumerate}
Група называје се директно неразложима, ако је неможно представити како директны продукт факторов $\ne\mE$.

Јасно јест: \emph{всака група с условјем минималности јест директны продукт конечно многих директно неразложимых груп.}\footnote{Како показали W. Krull в комутативном случају и О. Schmidt в обчем случају (ср. прим. 6), под предположеньем условја максималности и минималности представјенье јест једнозначно до изоморфизма. Понеже, без измены појма директној неразложимости, можно јест додати все внутришнье автоморфизмы к области операторов, достаточно јест жадати условја конечности за нормалне делителье или сушчествованье главного рјада.}

Ако се дотичне групы пишут адитивно, тогда појмы продукт и директны продукт преходет в суму и директну суму. Означенье: за суму груп $(\mA_1,\mA_2,\ldots)$, за директну суму $\mA_1+\mA_2+\cdots$.

Појем "директного пресека" далеј не буде употребјен, але наступне теоремы о связи меджу директным пресеком и директным продуктом показывајут, како теорију идеалов в наступных главах можно такоже формуловати с пресеками вместо сум; тако се установјаје связ со знаными комутативными теоријами.

Група $\mD$ называје се директным пресеком двух груп $\mR$ и $\mS$ в односу до $\mG$, ако
\[
\begin{array}{ll}
1.& \mR,\mS\text{ сут нормалне делительи в }\mG,\\
2.& \mR\cap\mS=\mD,\\
3.& \mR\mS=\mG.
\end{array}
\]
Подобно $\mD$ называје се директным пресеком $n$ груп $\mS_1,\ldots,\mS_n$, ако, положивши $\mR_i=\mS_1\cap\cdots\cap\mS_{i-1}\cap\mS_{i+1}\cap\cdots\cap\mS_n$, имајемо, же $\mD=\mR_i\cap\mS_i$ јест директны за всако $i$.

Из дефиниције следује, же $\mD$ мора быти нормалны делитель в $\mG$. Ако $\mD=\mS_1\cap\cdots\cap\mS_n$ јест директны и при хомоморфизму $\mG\sim\mG/\mD$ групы $\mG,\mS,\mD$ преходет в $\overline{\mG},\overline{\mS},\mE$, тогда $\mE=\overline{\mS}_1\cap\cdots\cap\overline{\mS}_n$ ставаје директны. Обратно, из всакого такого представјеньа $\mE=\overline{\mS}_1\cap\cdots\cap\overline{\mS}_n$ врачајемо се к директному представјеньу $\mD=\mS_1\cap\cdots\cap\mS_n$. Чрез ту одповеданье теоремы о директных пресеках при произвольном $\mD$ сводет се к специалному случају $\mD=\mE$.

В случају $\mD=\mE$ и двух факторов условја за директны продукт и директны пресек совпадајут. При $n$ факторах сушчествује наступна једино-једнозначна связ меджу директным продуктом и пресеком:

\emph{I. Ако $\mG=\mA_1\times\cdots\times\mA_n=\mA_i\times\mB_i$, тогда $\mE=\mB_1\cap\cdots\cap\mB_n=\mB_i\cap\mC_i$ јест директны и $\mC_i=\mA_i$.}

\emph{II. Ако $\mE=\mS_1\cap\cdots\cap\mS_n=\mR_i\cap\mS_i$ јест директны, тогда $\mG=\mR_1\times\cdots\times\mR_n=\mR_i\times\mT_i$ и $\mT_i=\mS_i$.}

\emph{Доказ I.} Нехај $\mB=\mB_1\cap\cdots\cap\mB_n=\mB_i\cap\mC_i$; тогда $\mC_i\supseteq\mA_i$. Покажемо, же $\mC_i\subseteq\mA_i$. Нехај, на примјер, $c$ јест в $\mC_1$; тогда $c$ јест в $\mB_2,\ldots,\mB_n$, зато компонентно представјенье по $\mA$ даваје:
\[
        c=a_1a_2\cdots a_n=a_1ea_3\cdots a_n=a_1a_2e\cdots a_n=\cdots=a_1\cdots a_{n-1}e.
\]
Понеже продукт всих $\mA_i$ јест директны, та представјеньа совпадајут; на всаком месту кроме првого једен раз стоји $e$, зато $c=a_1$, или $\mC_1\subseteq\mA_1$, то јест $\mC_1=\mA_1$, и одповедајуче $\mC_i=\mA_i$. Из того $\mB=\mB_i\cap\mC_i=\mB_i\cap\mA_i=\mE$; $\mB_i\mC_i=\mB_i\mA_i=\mG$, зато пресек јест директны.

\emph{Доказ II.} Нехај $\mH=\mR_1\cdots\mR_n=\mR_i\mT_i$; тогда $\mT_i\subseteq\mS_i$. Покажемо, же $\mS_i\subseteq\mT_i$. Нехај, на примјер, $s$ јест в $\mS_1$; тогда, посредством $\mG=\mS_i\times\mR_i$,
\[
        s=s_1e=s_2r_2=\cdots=s_nr_n.
\]
Сформујмо $t_1=er_2\cdots r_n=t_ir_i$; тогда, беручи до увагы поелементну комутативност $\mR_i$ и $\mS_i$, доставајемо
\[
        st_1^{-1}=s_it_i^{-1},
\]
то јест елемент всих $\mS_i$, и зато равны $e$. Зато $\mS_1\subseteq\mT_1$, или $\mT_1=\mS_1$, и одповедајуче $\mT_i=\mS_i$. Из того $\mH=\mR_i\mT_i=\mR_i\mS_i=\mG$ и $\mR_i\cap\mT_i=\mR_i\cap\mS_i=\mE$; следовательно $\mG$ јест директны продукт $\mR_i$.
\end{document}

